\section{Experiments}
\label{sec:experiments}

We want to verify whether the backbone of Section~\ref{sec:method},
trained with the contrastive loss of Section~\ref{sec:loss}, learns
information about the process that generated the data and, in
particular, whether the underlying parameters can be recovered from
the latents. We do this by training on a controlled synthetic
distribution where the generating parameters are known: stationary
autoregressive moving-average (ARMA) processes.\footnote{Code,
training scripts and experiment notebooks are available at
\url{https://github.com/jeremycochoy/contrastive-forecasting}.}

\subsection{ARMA training data and parameter sampling}
\label{sec:exp-data}

For each training sample we draw orders $p, q$ uniformly from
$\{1, \ldots, P\}$, then sample the $p$ AR coefficients and the $q$
MA coefficients independently from one of two distributions chosen by
a fair coin flip per sample: $\mathrm{Uniform}(-1, 1)$ or standard
normal $\mathcal{N}(0, 1)$. We then rescale each coefficient set so
that its $\ell_1$ norm does not exceed $0.95$: if
$\sum_i |\phi_i| > 0.95$, we multiply every $\phi_i$ by
$0.95 / \sum_i |\phi_i|$, and similarly for $\theta$. The bound
$\sum_i |\phi_i| \le 0.95$ is a sufficient stationarity condition for
the AR polynomial $1 - \sum_i \phi_i L^i$ (and an invertibility
condition for the MA polynomial $1 + \sum_i \theta_i L^i$). The
coefficient vectors are then zero-padded up to length $P$, so that
every sample has a fixed total of $2P$ ground-truth coefficients
regardless of the realised orders. We evaluate two configurations:
$P = 2$ (which we call \emph{2x2}, $4$ coefficients per sample) and
$P = 4$ (\emph{2x4}, $8$ coefficients per sample); 2x2 serves as a
smaller comparison point.

A series is generated by simulating $T = 4096$ steps from the sampled
ARMA process driven by i.i.d. standard normal noise. We use the
simplification $F = 1$ and $C = 1$ (each sample is a single univariate
series). The series is then patched into $T_p = T / W = 128$ patches
of width $W = 32$ following Section~\ref{sec:architecture}.

Training samples are drawn fresh at every step: every batch is a new
draw of orders, coefficients and noise, and there is no fixed dataset
of finite size.

\subsection{Loss}
\label{sec:exp-loss}

Section~\ref{sec:loss} defines three negative families. In the present
setting two of them simplify away.
\begin{itemize}
\item The cross-channel term $\mathcal{N}_{\mathrm{cc}}$ vanishes
  because $C = 1$: there is no other channel to draw negatives from.
\item The cross-temporal term $\mathcal{N}_{\mathrm{ct}}$ is dropped by
  design. We want to learn the underlying parameters of the
  generating ARMA process, not the specific realised samples;
  successive samples in the same series are produced by the same
  parameters and the past directly influences the future, so a
  cross-temporal pair is two views of the same process, and pushing
  them apart in latent space would conflict with our goal.
\end{itemize}
Only the cross-batch term $\mathcal{N}_{\mathrm{cb}}$ remains.
Dropping the channel index, the contribution of an anchor $(b, t)$ to
the loss is
\begin{equation}
  \mathcal{L}_{b, t}
  \;=\; -\log
  \frac{\mathcal{P}_+(b, t)}
  {\mathcal{N}_{\mathrm{cb}}(t)},
\end{equation}
with
\[
  \mathcal{P}_+(b, t)
  \;=\; \sigma_\tau\!\big(f[b, t],\, h[b, t+1]\big),
  \qquad
  \mathcal{N}_{\mathrm{cb}}(t)
  \;=\; \sum_{b \neq b'}
    \sigma_\tau\!\big(h[b, t+1],\, f[b', t]\big).
\]
We use $\tau = 0.07$.

\subsection{Backbone}
\label{sec:exp-backbone}

We use the backbone of Section~\ref{sec:method} with $L = 12$
transformer layers and hidden width $H = 1024$. The encoder is the
bidirectional GRU described in Section~\ref{sec:encoder}. Training
runs for $2$M steps with AdamW, learning rate $7 \times 10^{-5}$,
batch size $8$. The architecture-search history that motivates these
choices is reported in Appendix~\ref{sec:appendix-backbone-search}.

We monitor backbone training with the \emph{contrastive gap}, defined
as
\[
  \mathrm{gap} \;=\; \mathbb{E}_{b, t} \big[
    \langle f[b, t],\, h[b, t+1] \rangle
    \;-\;
    \langle f[b, t],\, h[b, t] \rangle
  \big],
\]
the difference between the mean forecast-future and forecast-past
cosine similarities on a held-out batch. A higher gap means the
forecast latent is closer to the future patch at $t + 1$ than to the
past patch at $t$. The gap is a diagnostic, not optimised directly;
it tracks well with downstream recovery quality and is the metric we
use to compare backbone variants in
Appendix~\ref{sec:appendix-backbone-search}.

\subsection{Recovery head}
\label{sec:exp-head}

To probe whether the backbone latents encode the ARMA parameters, we
freeze the backbone after training and fit a small \emph{recovery
head} on top of it. The head receives the per-patch latent
$h[b, t] \in \mathbb{R}^H$ produced by the frozen backbone and returns
a per-patch coefficient prediction
$\hat{\theta}[b, t] \in \mathbb{R}^{2P}$.

Internally the head has four parts: an input projection (linear,
GELU, layer norm) mapping $\mathbb{R}^H \to \mathbb{R}^{128}$, a
bidirectional GRU with hidden size $128$ and $2$ layers, a two-layer
GELU MLP applied to the GRU output, and two parallel linear heads
(one for AR, one for MA), each followed by $\tanh$ to constrain
outputs to $(-1, 1)$ in line with the sampling range. The full head
has roughly $676$k parameters. The single coefficient
prediction for a series is the mean of the per-patch predictions:
\[
  \hat{\theta}[b]
  \;=\;
  \frac{1}{T_p} \sum_{t = 1}^{T_p} \hat{\theta}[b, t].
\]
The head is trained with mean squared error against the ground-truth
coefficients of each sample. The architecture-search history that led
to this head is reported in Appendix~\ref{sec:appendix-head-search}.

We evaluate the head with the \emph{improvement ratio}
\begin{equation}
  r \;=\; \frac{\mathrm{MSE}_0}{\mathrm{MSE}_{\hat{\theta}}},
  \label{eq:improvement-ratio}
\end{equation}
where $\mathrm{MSE}_0$ is the MSE of the all-zero baseline that
predicts $\hat{\theta} = 0$ for every coefficient. This baseline is
non-trivial: since order padding leaves a fraction of the ground-truth
coefficients at exactly zero, the all-zero predictor is already
perfect on those positions and so achieves a fairly low MSE. The head
must do better still in order to reach $r > 1$.

\subsection{Random-walk correlation experiment}
\label{sec:exp-corr}

We run a second probing experiment in which the targets are pairwise
correlations of a multivariate random walk.  For each training sample we
draw a $4\times4$ correlation matrix $\Sigma$ by sampling the six
off-diagonal entries i.i.d.\ from $\mathrm{Uniform}(0,1)$ and rejecting
the matrix if it is not positive semi-definite (acceptance rate
${\approx}\,43\%$, vectorised over an oversampled candidate batch).  A
series of $C=4$ channels of length $T=4096$ is generated by cumulating
i.i.d.\ increments drawn from $\mathcal{N}(0,\Sigma)$; each channel is
z-scored within the sample.  The six upper-triangular off-diagonal entries
of $\Sigma$ are the recovery targets.

The backbone uses a \emph{joint-channel} design: per-channel GRU encodings
at each timestep are concatenated and projected to $\mathbb{R}^H$ before
the transformer, collapsing the channel dimension to $1$.  The transformer
therefore sees a single sequence $[B,T,H]$ per sample.  As a consequence,
$\mathcal{N}_{\mathrm{cc}}$ vanishes ($C=1$ from the loss's perspective)
and $\mathcal{N}_{\mathrm{ct}}$ is dropped for the same reason as in the
ARMA case (successive patches share the same $\Sigma$).  Only
$\mathcal{N}_{\mathrm{cb}}$ is active, giving the same loss form as
Eq.~\eqref{eq:loss-anchor}.  The backbone uses $L=12$, $H=1024$, batch
size $16$, and is trained for $195\,000$ steps with AdamW at
$7\times10^{-5}$.  Peak contrastive gap: $0.506$ at step $127\,000$.

The recovery head mirrors the architecture of
Section~\ref{sec:exp-head}: input projection, bidirectional GRU ($128$
units, $2$ layers), mean-pool over time, two-layer GELU MLP, and a
linear output of dimension $6$ followed by $\mathrm{clamp}(0,1)$.  It
is trained with MSE against the six ground-truth correlations, with the
backbone frozen.
